% 翻译状态：专题标题已译；正文待继续翻译与校对。
% Chapter 3, Topic _Linear Algebra_ Jim Hefferon
%  http://joshua.smcvt.edu/linearalgebra
%  2001-Jun-12
\topic{最佳拟合直线}
\index{least squares|(}
\index{line of best fit|(}
\index{linear equation!inconsistent systems}
\textit{本专题需要使用本章相关小节中的公式，
        正交投影到直线和投影到子空间。} 

科学家经常遇到无解的方程组，却仍必须给出答案。更准确地说，他们必须找到最佳答案。
例如，下面是抛掷硬币的结果，其中包括若干中间数据。
\begin{center} \small
  \begin{tabular}{r|ccc}
     \textit{number of flips}  &30  &60  &90   \\
     \hline
     \textit{number of heads}  &16  &34  &51
  \end{tabular}
\end{center}
由于实验具有随机性，正面数与抛掷次数之比会围绕硬币长期的 50-50 比例波动。因此，这类实验对应的方程组很可能无解，本例正是如此。 
\begin{equation*}
  \begin{linsys}{1}
    30m  &=  &16    \\
    60m  &=  &34    \\
    90m  &=  &51
  \end{linsys}
\end{equation*}
也就是说，我们收集的数据向量不在理想情况下应属于的子空间中。
\begin{equation*}
  \colvec{16 \\ 34 \\ 51}\not\in
    \set{ m\colvec{30 \\ 60 \\ 90}  \suchthat m\in\Re}
\end{equation*}
但我们仍需采取措施，于是寻找最接近成立的 $m$。将数据向量正交投影到直线子空间，得到最佳估计，即子空间中距离数据向量最近的向量。
\begin{equation*}
  \frac{ \colvec[r]{16 \\ 34 \\ 51}\dotprod\colvec[r]{30 \\ 60 \\ 90} }{
           \colvec[r]{30 \\ 60 \\ 90}\dotprod\colvec[r]{30 \\ 60 \\ 90} }
     \cdot\colvec[r]{30 \\ 60 \\ 90}
  =\frac{7110}{12600}\cdot \colvec[r]{30 \\ 60 \\ 90}
\end{equation*}
估计值（\( m=7110/12600\approx 0.56 \)）略大于一半，但差别不大，因此这枚硬币大概是公平的。 

斜率为 \( m\approx 0.56 \) 的直线称为该数据的\definend{最佳拟合直线}\index{line!best fit}%
\index{best fit line} 
用于这些数据。
\begin{center}  \small
  \includegraphics{map/mp/ch3.42}
\end{center}
使给定向量与作为等式右端向量之间的距离最小，就使这些竖直长度之和最小，因此称该直线由\definend{最小二乘拟合}得到。\index{least squares}
\begin{center}  \small
  \includegraphics{map/mp/ch3.43}
\end{center}
为使长度更清晰，图中将竖直方向的比例放大了十倍。

上式中的直线必须经过 \( (0,0) \)，因为我们把它视为斜率等于该硬币正面比例的直线。我们也能处理直线不经过原点的情形。

下面给出男子一英里赛世界纪录用时的变化 \cite{Oakley}。
20 世纪初，许多人想知道这项纪录何时（或是否）会降到 4 分钟以内。下表列出该世纪前半叶每十年元旦时保持的纪录。
\begin{center} \small
  \begin{tabular}{r|ccccccccc}
    \textit{year} &$1870$ &$1880$ &$1890$  &$1900$  
        &$1910$  &$1920$  &$1930$ &$1940$ &$1950$   \\
    \hline
    \textit{secs}  &$268.8$  &$264.5$  &$258.4$  &$255.6$  
        &$255.6$  &$252.6$  &$250.4$ &$246.4$ &$241.4$
  \end{tabular}  
\end{center}
我们可据此预测约在 1950 年达到 $240$ 秒的日期，再与实际日期比较。和硬币数据一样，这些数值不在一条完美直线上；也就是说，该方程组对斜率和截距没有精确解。
\begin{equation*}
  \begin{linsys}{2}
    b &+  &1870m  &=  &268.8 \\ 
    b &+  &1880m  &=  &264.5 \\ 
      &   &       &\vdots     \\
    b &+  &1950m  &=  &241.4  
  \end{linsys}
\end{equation*}
我们使用正交投影求最佳近似。

（\textit{数据说明。}仅取每十年开始时的记录可减少数据录入量，并在一定程度上平滑数据；与录入所有日期和纪录所得结果大致相同。不同标准机构采用的时间序列有所不同，这里采用的数据来自 \cite{WikipediaMensMile}。图从 1870 年开始，是因为当时曾有“职业”和“业余”两类纪录；后来前一类不再更新，因此我们采用后一类。）

写出该线性方程组的系数矩阵和常数向量（世界纪录用时）。
\begin{equation*}
  A=
  \begin{mat}
    1  &1870  \\
    1  &1880  \\
    \vdots  &\vdots \\
    1  &1950  
  \end{mat}
  \qquad
  \vec{v}=\colvec{268.8 \\ 264.5 \\ \vdots \\ 241.4}
\end{equation*}
“投影到子空间”小节的结论
给出了系数
$b$ and $m$ that make the linear combination
of $A$'s columns as close as possible to $\vec{v}$. 
这些系数就是向量 $(\trans{A}A)^{-1}\trans{A}\cdot\vec{v}$ 的各个分量。

\textit{Sage} can do the computation for us.
\begin{lstlisting}
sage: year = [1870, 1880, 1890, 1900, 1910, 1920, 1930, 1940, 1950]
sage: secs = [268.8, 264.5, 258.4, 255.6, 255.6, 252.6, 250.4, 246.4, 241.4]
sage: var('a, b, t')
(a, b, t)
sage: model(t) = a*t+b
sage: data = zip(year, secs)
sage: fit = find_fit(data, model, solution_dict=True)
sage: model.subs(fit)
t |--> -0.3048333333333295*t + 837.0872222222147
sage: g=points(data)+plot(model.subs(fit),(t,1860,1960),color='red',
....:                     figsize=3,fontsize=7,typeset='latex')
sage: g.save("four_minute_mile.pdf")
sage: g
\end{lstlisting}

\begin{center}  
  \includegraphics[width=0.5\textwidth]{map/pix/four_minute_mile.pdf}
\end{center}
这些数据形成的直线拟合效果出人意料地好。由斜率和截距预测日期为 $1958.73$；罗杰·班尼斯特的纪录实际日期是 1954-05-06。

最后一个例子比较 2002 年美国职业棒球大联盟各队的薪资与获胜场次。该年奥克兰运动家队据电影 \textit{Moneyball} 所述，运用数学方法在预算内优化上场球员。（薪资单位为百万美元，获胜场次以 162 场比赛为总数。） 

计算仍使用 \textit{Sage}。 
% http://ask.sagemath.org/question/768/fitting-a-curve-to-a-straight-line?answer=1271#1271
% by Joaquim Puig on 2011-Sep-27; gathered on 2013-Oct-23 JH
\begin{lstlisting}
sage: sal = [40, 40, 39, 42, 45, 42, 62, 34, 41, 57, 58, 63, 47, 75, 57, 78, 80, 50, 60, 93,
....:        77, 55, 95, 103, 79, 76, 108, 126, 95, 106]
sage: wins = [103, 94, 83, 79, 78, 72, 99, 55, 66, 81, 80, 84, 62, 97, 73, 95, 93, 56, 67,
....:        101, 78, 55, 92, 98, 74, 67, 93, 103, 75, 72]
sage: var('a, b, t')
(a, b, t)
sage: model(t) = a*t+b
sage: data = zip(sal,wins)
sage: fit = find_fit(data, model, solution_dict=True)
sage: model.subs(fit)
t |--> 0.2634981251436269*t + 63.06477642781477
sage: p = points(data,size=25)+plot(model.subs(fit),(t,30,130),color='red',typeset='latex')
sage: p.save('moneyball.pdf')
\end{lstlisting}

图如下。左上方以较低薪资取得较多胜场的球队是奥克兰 A 队。

\begin{center}  \small
  \includegraphics[width=0.5\textwidth]{map/pix/moneyball.pdf}
\end{center}

仅凭目测判断这条直线容易出错。因此，方程能确定最佳拟合中的“最佳”。此外，模型方程大致表明：球队老板每多花一百万美元，预计可多获得 $1/4$ 场胜利（当然，这一预期并不十分可靠）。


% For example, the different denominations of US\ money have different average 
% times in circulation.\cite{FedReserve}
% (The $\$2$~bill is a special case because 
% Americans mistakenly believe that 
% it is collectible and do not circulate these bills.)
% \begin{lstlisting}
% sage: x = [1, 5, 10, 20, 50, 100]
% sage: y = [5.9, 4.9, 4.2, 7.7, 3.7, 15.0]
% sage: var('a, b, t') 
% (a, b, t)
% sage: model(t) = a*t+b
% sage: data = zip(x,y)
% sage: fit = find_fit(data, model, solution_dict=True)
% sage: model.subs(fit)                                
% t |--> 0.08650137692301652*t + 4.2184573359365265
% sage: plot(model.subs(fit),(t,0,100))+points(data,size=20,color='red')
% \end{lstlisting}  
% How long should a $\$25$~bill last?

% \begin{center}  \small
%   \includegraphics[width=0.5\textwidth]{map/pix/denom.png}
% \end{center}


% \begin{center}  \small
%   \begin{tabular}{r|cccccc}
%      \textit{denomination}          
%         &$1$   &$5$    &$10$     &$20$     &$50$     &$100$ \\
%       \hline
%      \textit{average life (mos)}  
%         &$22.0$ &$15.9$  &$18.3$   &$24.3$   &$55.4$   &$88.8$  \\
%   \end{tabular}
% \end{center}
% The data plot below looks roughly linear.
% It isn't a perfect line, i.e.,
% the linear system with equations $b+1m=1.5$, \ldots, $b+100m=20$ has no 
% solution, but we can again use orthogonal projection 
% to find a best approximation.
% 考虑 matrix of coefficients of that linear system and also its
% vector of constants, the experimentally-determined values.
% \begin{equation*}
%   A=
%   \begin{mat}[r]
%     1  &1  \\
%     1  &5  \\
%     1  &10  \\
%     1  &20  \\
%     1  &50  \\
%     1  &100  
%   \end{mat}
%   \qquad
%   \vec{v}=\colvec[r]{22.0 \\ 15.9 \\ 18.3 \\ 24.3 \\ 55.4 \\ 88.8}
% \end{equation*}
% “投影到子空间”小节的结论says
% that coefficients $b$ and $m$ so that the linear combination
% of the columns of $A$ is as close as possible to the vector $\vec{v}$ 
% are the entries of $(\trans{A}A)^{-1}\trans{A}\cdot\vec{v}$.
% Some calculation
% gives an intercept of $b=14.16$ and a slope
% of $m= 0.75$.
% \begin{center}  \small
%   \includegraphics{map/mp/ch3.44}
% \end{center}
% Plugging $x=25$ into the equation of the line shows that such a 
% bill should last about two and three-quarters years.


%
%
%We want to project when a 3:40~mile %\hrsmins{3}{40} 
%will be run.
%\begin{center}
%  \begin{tabular}{r|ccccccc}
%    \textit{year} &1870  &1880  &1890  &1900  &1910  &1920  &1930   \\
%    \hline
%    \textit{secs}  &$268.8$  &$264.5$  &$258.4$  &$255.6$  
%        &$255.6$  &$252.6$  &$250.4$
%  \end{tabular}        \\[2ex]
%  \begin{tabular}{|ccccccc}
%      1940  &1950  &1960  &1970  &1980  &1990 &2000                  \\
%      \hline
%      $246.4$  &$241.4$  &$234.5$  &$231.1$  &$229.0$  &$226.3$ &$223.1$
%  \end{tabular}
%\end{center}
%We can see below that the data is surprisingly linear.
%With this input
%\begin{equation*}
%  A= 
%  \begin{mat}
%    1 & 1860 \\
%    1 & 1870 \\
%%    1 & 1880 \\
%%    1 & 1890 \\ 
%%    1 & 1900 \\ 
%%    1 & 1910 \\
%%    1 & 1920 \\
%%    1 & 1930 \\
%%    1 & 1940 \\
%%    1 & 1950 \\
%%    1 & 1960 \\
%%    1 & 1970 \\
%%    1 & 1980 \\
%    \vdots &\vdots  \\
%    1 &1990  \\
%    1 &2000
%  \end{mat}
%  \qquad
%  \vec{v} = \colvec{280.0 \\
%                    268.8 \\
%%                    264.5 \\
%%                    258.4 \\
%%                    255.6 \\
%%                    255.6 \\
%%                    252.6 \\
%%                    250.4 \\
%%                    246.4 \\
%%                    241.4 \\
%%                    234.5 \\
%%                    231.1 \\
%%                    229.0 \\
%                    \vdots \\
%                    226.3 \\
%                    223.1}
%\end{equation*}
%%MAPLE gives $b=970.68$ and $m=-0.37$  (rounded to two places).
%the Python program at this Topic's end gives 
%%$\text{slope}= -0.351428571429$
%%and $\text{intercept}= 925.534065934$ 
%$\text{slope}= -0.35$
%and $\text{intercept}= 925.53$ 
%(rounded to two places; the original data 
%is good to only about a quarter of a second 
%since much of it was hand-timed).
%\begin{center}  \small
%  \includegraphics{map/mp/ch3.45}
%\end{center}
%When will a \( 220 \) second mile be run?
%Solving the equation of the line of best fit 
%gives an estimate of the year \( 2008 \).
%
%This example is amusing, but serves as a caution \Dash obviously the 
%linearity of the data will break down someday (as indeed it does 
%prior to 1860).



\begin{exercises}
  \item[]\emph{这里的计算最好使用计算机完成。部分习题需要从互联网获取数据。}
%   \begin{Filesave}{bookans}  % write the records to the answer file.
%     \textit{Data on the progression of the world's records 
%        (taken from the \textit{Runner's World} web site)       
%        is below.}
%     \begin{figure}
%       \rotatebox[origin=lb]{90}{\hbox{\scriptsize
%         \begin{tabular}[t]{|l|ll|}
% %        \hline
% %          \multicolumn{3}{|c|}{\textit{Amateurs}}    \\
%         \hline
%           \multicolumn{3}{|c|}{\textit{Progression of Men's Mile Record}}  \\
%         \hline
%         \textit{time}  &\textit{name}  &\textit{date}  \\
%          \hline  
%            $\text{4:52}.0$      &Cadet Marshall (GBR)       &02Sep52 \\
%            $\text{4:45}.0$      &Thomas Finch (GBR)         &03Nov58 \\
%            %$\text{4:45}.0$      &St.~Vincent Hammick (GBR)  &15Nov58 \\
%            $\text{4:40}.0$      &Gerald Surman (GBR)        &24Nov59 \\
%            $\text{4:33}.0$      &George Farran (IRL)        &23May62 \\
%            $\text{4:29}\;3/5$   &Walter Chinnery (GBR)      &10Mar68 \\
%            $\text{4:28}\;4/5$   &William Gibbs (GBR)        &03Apr68 \\ 
%            $\text{4:28}\;3/5$   &Charles Gunton (GBR)       &31Mar73 \\
%            $\text{4:26}.0$      &Walter Slade (GBR)         &30May74 \\
%            $\text{4:24}\;1/2$   &Walter Slade (GBR)         &19Jun75 \\
%            $\text{4:23}\;1/5$   &Walter George (GBR)        &16Aug80 \\
%            $\text{4:19}\;2/5$   &Walter George (GBR)        &03Jun82 \\
%            $\text{4:18}\;2/5$   &Walter George (GBR)        &21Jun84 \\
%            $\text{4:17}\;4/5$   &Thomas Conneff (USA)       &26Aug93 \\
%            $\text{4:17}.0$      &Fred Bacon (GBR)           &06Jul95 \\
%            $\text{4:15}\;3/5$   &Thomas Conneff (USA)       &28Aug95 \\
%            $\text{4:15}\;2/5$   &John Paul Jones (USA)      &27May11 \\
%            $\text{4:14}.4$      &John Paul Jones (USA)      &31May13 \\
%            $\text{4:12}.6$      &Norman Taber (USA)         &16Jul15 \\
%            $\text{4:10}.4$      &Paavo Nurmi (FIN)          &23Aug23 \\
%            $\text{4:09}\;1/5$   &Jules Ladoumegue (FRA)     &04Oct31 \\
%            $\text{4:07}.6$      &Jack Lovelock (NZL)        &15Jul33 \\
%            $\text{4:06}.8$      &Glenn Cunningham (USA)     &16Jun34 \\
%            $\text{4:06}.4$      &Sydney Wooderson (GBR)     &28Aug37 \\
%            $\text{4:06}.2$      &Gunder Hagg (SWE)          &01Jul42 \\
%            %$\text{4:06}.2$      &Arne Andersson (SWE)       &10Jul42 \\
%            $\text{4:04}.6$      &Gunder Hagg (SWE)          &04Sep42 \\
%            $\text{4:02}.6$      &Arne Andersson (SWE)       &01Jul43 \\
%            $\text{4:01}.6$      &Arne Andersson (SWE)       &18Jul44 \\
%            $\text{4:01}.4$      &Gunder Hagg (SWE)          &17Jul45 \\
%            $\text{3:59}.4$      &Roger Bannister (GBR)      &06May54 \\
%            $\text{3:58}.0$      &John Landy (AUS)           &21Jun54 \\
%            $\text{3:57}.2$      &Derek Ibbotson (GBR)       &19Jul57 \\
%            $\text{3:54}.5$      &Herb Elliott (AUS)         &06Aug58 \\
%            $\text{3:54}.4$      &Peter Snell (NZL)          &27Jan62 \\ 
%            $\text{3:54}.1$      &Peter Snell (NZL)          &17Nov64 \\
%            $\text{3:53}.6$      &Michel Jazy (FRA)          &09Jun65 \\
%            $\text{3:51}.3$      &Jim Ryun (USA)             &17Jul66 \\
%            $\text{3:51}.1$      &Jim Ryun (USA)             &23Jun67 \\
%            $\text{3:51}.0$      &Filbert Bayi (TAN)         &17May75 \\
%            $\text{3:49}.4$      &John Walker (NZL)          &12Aug75 \\
%            $\text{3:49}.0$      &Sebastian Coe (GBR)        &17Jul79 \\
%            $\text{3:48}.8$      &Steve Ovett (GBR)          &01Jul80 \\
%            $\text{3:48}.53$     &Sebastian Coe (GBR)        &19Aug81 \\
%            $\text{3:48}.40$     &Steve Ovett (GBR)          &26Aug81 \\
%            $\text{3:47}.33$     &Sebastian Coe (GBR)        &28Aug81 \\
%            $\text{3:46}.32$     &Steve Cram (GBR)           &27Jul85 \\
%            $\text{3:44}.39$     &Noureddine Morceli (ALG)   &05Sep93 \\
% 	   $\text{3:43}.13$     &Hicham el Guerrouj (MOR)   &07Jul99 \\
%       \hline
%       \end{tabular}
%       \   %
%         \begin{tabular}[t]{|l|ll|}
%         \hline
%           \multicolumn{3}{|c|}{\textit{Progression of Men's $1500$~Meter Record}}  \\
%         \hline
%         \textit{time}  &\textit{name}  &\textit{date}  \\
%          \hline  
%          $\text{4:09}.0$   &John Bray (USA)        &30May00  \\
%          $\text{4:06}.2$   &Charles Bennett (GBR)  &15Jul00  \\
%          $\text{4:05}.4$   &James Lightbody (USA)  &03Sep04  \\
%          $\text{3:59}.8$   &Harold Wilson (GBR)    &30May08  \\
%          $\text{3:59}.2$   &Abel Kiviat (USA)      &26May12  \\
%          $\text{3:56}.8$   &Abel Kiviat (USA)      &02Jun12  \\
%          $\text{3:55}.8$   &Abel Kiviat (USA)      &08Jun12  \\
%          $\text{3:55}.0$   &Norman Taber (USA)     &16Jul15  \\
%          $\text{3:54}.7$   &John Zander (SWE)      &05Aug17  \\
%          $\text{3:53}.0$   &Paavo Nurmi (FIN)      &23Aug23  \\
%          $\text{3:52}.6$   &Paavo Nurmi (FIN)      &19Jun24  \\
%          $\text{3:51}.0$   &Otto Peltzer (GER)     &11Sep26  \\
%          $\text{3:49}.2$   &Jules Ladoumegue (FRA) &05Oct30  \\
%          %$\text{3:49}.2$   &Luigi Beccali (ITA)    &09Sep33  \\
%          $\text{3:49}.0$   &Luigi Beccali (ITA)    &17Sep33  \\
%          $\text{3:48}.8$   &William Bonthron (USA) &30Jun34  \\
%          $\text{3:47}.8$   &Jack Lovelock (NZL)    &06Aug36  \\
%          $\text{3:47}.6$   &Gunder Hagg (SWE)      &10Aug41  \\
%          $\text{3:45}.8$   &Gunder Hagg (SWE)      &17Jul42  \\
%          $\text{3:45}.0$   &Arne Andersson (SWE)   &17Aug43  \\
%          $\text{3:43}.0$   &Gunder Hagg (SWE)      &07Jul44  \\
%          %$\text{3:43}.0$   &Lennart Strand (SWE)   &15Jul47  \\
%          %$\text{3:43}.0$   &Werner Lueg (FRG)      &29Jun52  \\
%          %$\text{3:43}.0$   &Roger Bannister (GBR)  &06May54  \\
%          $\text{3:42}.8$   &Wes Santee (USA)       &04Jun54  \\
%          $\text{3:41}.8$   &John Landy (AUS)       &21Jun54  \\
%          $\text{3:40}.8$   &Sandor Iharos (HUN)    &28Jul55  \\
%          %$\text{3:40}.8$   &Laszlo Tabori (HUN)    &06Sep55  \\
%          %$\text{3:40}.8$   &Gunnar Nielsen (DEN)   &06Sep55  \\
%          $\text{3:40}.6$   &Istvan Rozsavolgyi (HUN) &03Aug56 \\
%          $\text{3:40}.2$   &Olavi Salsola (FIN)    &11Jul57  \\
%          %$\text{3:40}.2$   &Olavi Salonen (FIN)    &11Jul57  \\
%          $\text{3:38}.1$   &Stanislav Jungwirth (CZE) &12Jul57  \\
%          $\text{3:36}.0$   &Herb Elliott (AUS)     &28Aug58  \\
%          $\text{3:35}.6$   &Herb Elliott (AUS)     &06Sep60  \\
%          $\text{3:33}.1$   &Jim Ryun (USA)         &08Jul67  \\
%          $\text{3:32}.2$   &Filbert Bayi (TAN)     &02Feb74  \\
%          $\text{3:32}.1$   &Sebastian Coe (GBR)    &15Aug79  \\
%          %$\text{3:32}.1$   &Steve Ovett (GBR)      &15Jul80  \\  
%          $\text{3:31}.36$  &Steve Ovett (GBR)      &27Aug80  \\
%          $\text{3:31}.24$  &Sydney Maree (usa)     &28Aug83  \\
%          $\text{3:30}.77$  &Steve Ovett (GBR)      &04Sep83  \\
%          $\text{3:29}.67$  &Steve Cram (GBR)       &16Jul85  \\
%          $\text{3:29}.46$  &Said Aouita (MOR)      &23Aug85  \\ 
%          $\text{3:28}.86$  &Noureddine Morceli (ALG)  &06Sep92  \\
%          $\text{3:27}.37$  &Noureddine Morceli (ALG)  &12Jul95  \\
%          $\text{3:26}.00$  &Hicham el Guerrouj (MOR)  &14Jul98  \\
%      \hline
%      \end{tabular}
%       \   %
%         \begin{tabular}[t]{|l|ll|}
%         \hline
%           \multicolumn{3}{|c|}{\textit{Progression of Women's Mile Record}}  \\
%         \hline
%         \textit{time}  &\textit{name}  &\textit{date}  \\
%          \hline  
%            $\text{6:13}.2$    &Elizabeth Atkinson (GBR)    &24Jun21 \\
%            $\text{5:27}.5$    &Ruth Christmas (GBR)        &20Aug32 \\
%            $\text{5:24}.0$    &Gladys Lunn (GBR)           &01Jun36 \\
%            $\text{5:23}.0$    &Gladys Lunn (GBR)           &18Jul36 \\
%            $\text{5:20}.8$    &Gladys Lunn (GBR)           &08May37 \\
%            $\text{5:17}.0$    &Gladys Lunn (GBR)           &07Aug37 \\
%            $\text{5:15}.3$    &Evelyne Forster (GBR)       &22Jul39 \\
%            $\text{5:11}.0$    &Anne Oliver (GBR)           &14Jun52 \\
%            $\text{5:09}.8$    &Enid Harding (GBR)          &04Jul53 \\
%            $\text{5:08}.0$    &Anne Oliver (GBR)           &12Sep53 \\
%            $\text{5:02}.6$    &Diane Leather (GBR)         &30Sep53 \\
%            $\text{5:00}.3$    &Edith Treybal (ROM)         &01Nov53 \\
%            $\text{5:00}.2$    &Diane Leather (GBR)         &26May54 \\
%            $\text{4:59}.6$    &Diane Leather (GBR)         &29May54 \\
%            $\text{4:50}.8$    &Diane Leather (GBR)         &24May55 \\
%            $\text{4:45}.0$    &Diane Leather (GBR)         &21Sep55 \\
%            $\text{4:41}.4$    &Marise Chamberlain (NZL)    &08Dec62 \\
%            $\text{4:39}.2$    &Anne Smith (GBR)            &13May67 \\
%            $\text{4:37}.0$    &Anne Smith (GBR)            &03Jun67 \\
%            $\text{4:36}.8$    &Maria Gommers (HOL)         &14Jun69 \\
%            $\text{4:35}.3$    &Ellen Tittel (FRG)          &20Aug71 \\
%            $\text{4:34}.9$    &Glenda Reiser (CAN)         &07Jul73 \\
%            $\text{4:29}.5$    &Paola Pigni-Cacchi (ITA)    &08Aug73 \\
%            $\text{4:23}.8$    &Natalia Marasescu (ROM)     &21May77 \\
%            $\text{4:22}.1$    &Natalia Marasescu (ROM)     &27Jan79 \\
%            $\text{4:21}.7$    &Mary Decker (USA)           &26Jan80 \\
%            $\text{4:20}.89$   &Lyudmila Veselkova (SOV)    &12Sep81 \\
%            $\text{4:18}.08$   &Mary Decker-Tabb (USA)      &09Jul82 \\
%            $\text{4:17}.44$   &Maricica Puica (ROM)        &16Sep82 \\
%            $\text{4:15}.8$    &Natalya Artyomova (SOV)     &05Aug84 \\
%            $\text{4:16}.71$   &Mary Decker-Slaney (USA)    &21Aug85 \\
%            $\text{4:15}.61$   &Paula Ivan (ROM)            &10Jul89 \\
%            $\text{4:12}.56$   &Svetlana Masterkova (RUS)   &14Aug96 \\
%        \hline
%        \end{tabular}
%       }}  %end the rotatebox
%     \end{figure}
%   \end{Filesave}  
  \item
    使用最小二乘法判断本实验中的硬币是否公平。
    \begin{center}
      \begin{tabular}{r|ccccc}
         \textit{flips}  &8  &16  &24  &32  &40 \\
         \hline
         \textit{heads}  &4  &9   &13  &17  &20
      \end{tabular}
    \end{center}
    \begin{answer}
      与前面的第一个例子一样，我们要寻找
      最佳 $m$ 来“求解”该方程组。 
      \begin{equation*}
        \begin{linsys}{5}
          8m  &=  &4  \\
          16m &=  &9  \\
          24m &=  &13 \\
          32m &=  &17 \\
          40m &=  &20
        \end{linsys}
      \end{equation*}
      投影到线性子空间得到
      \begin{equation*}
        \frac{\colvec[r]{4 \\ 9 \\ 13 \\ 17 \\20}
          \dotprod
          \colvec[r]{8  \\ 16 \\ 24 \\ 32 \\ 40}}{%
          \colvec[r]{8  \\ 16 \\ 24 \\ 32 \\ 40}
          \dotprod
          \colvec[r]{8  \\ 16 \\ 24 \\ 32 \\ 40}}
        \cdot
        \colvec[r]{8  \\ 16 \\ 24 \\ 32 \\ 40}
        =
        \frac{1832}{3520}
        \cdot
        \colvec[r]{8  \\ 16 \\ 24 \\ 32 \\ 40}
      \end{equation*}
      因此最佳拟合直线斜率约为 $0.52$。
      \begin{center}  \small
        \includegraphics{map/mp/ch3.85}
      \end{center}
    \end{answer}
  \item 
     对男子一英里纪录，与其列出许多
     纪录及其确切日期，不如通过
     定期抽样对数据作适度“平滑”。
     完成更长的计算并比较结论。
     \begin{answer}
     使用以下输入
     \begin{equation*}
       A =  
       \begin{mat}
           1 & 1852.71 \\
           1 & 1858.88 \\
          \vdots  &\vdots      \\
           1 & 1985.54 \\
           1 & 1993.71
        \end{mat}
        \qquad
        b = \colvec{ 292.0 \\ 
                     285.0 \\ 
                    \vdots  \\
                     226.32 \\
                     224.39}
     \end{equation*}
     (the dates have been rounded to months, e.g., for a September record,
     小数 $.71\approx (8.5/12)$），Maple 给出
     截距为 $b=994.8276974$，斜率为
     $m=-0.3871993827$.    
     \begin{center}  \small
        \includegraphics{map/mp/ch3.86}
     \end{center}
     \end{answer}
  \item 
    求出男子 $1500$ 米赛跑的最佳拟合直线。
    比较其斜率与男子一英里赛的斜率。
    (The distances are close; a mile is about $1609$~meters.)
    \begin{answer}
   使用以下输入（年份以 $1900$ 年为零点）
   \begin{equation*}
     A := 
     \begin{mat}
        1 &   .38   \\
        1 &   .54   \\
       \vdots \vdots \\ 
        1 & 92.71 \\
        1 & 95.54
     \end{mat}
     \qquad      
     b = \colvec{ 249.0 \\
                  246.2 \\
                 \vdots \\
                  208.86 \\
                  207.37}
   \end{equation*}
     (the dates have been rounded to months, e.g., for a September record,
     使用小数 $.71\approx (8.5/12)$），
     Maple 给出截距 $b=243.1590327$，斜率
   $m=-0.401647703$.
   本专题正文中男子一英里赛的斜率与此
   与此非常接近。
   \begin{center}  \small
     \includegraphics{map/mp/ch3.87}
   \end{center}
   \end{answer}
  \item 求出女子一英里纪录的最佳拟合直线。
    \begin{answer}
    使用以下输入（年份以 $1900$ 年为零点）
    \begin{equation*}
      A =  
      \begin{mat}
         1 & 21.46 \\
         1 & 32.63 \\
        \vdots  &\vdots     \\
         1 & 89.54  \\
         1 & 96.63
      \end{mat}
      \qquad
      b = 
      \colvec{ 373.2 \\
               327.5 \\
              \vdots  \\
               255.61 \\
               252.56}      
    \end{equation*}
     (the dates have been rounded to months, e.g., for a September record,
     使用小数 $.71\approx (8.5/12)$），
     MAPLE 给出截距 $b=378.7114894$，斜率
    $m=-1.445753225$.
    \begin{center}  \small
      \includegraphics{map/mp/ch3.88}
    \end{center}
   \end{answer}   
  \item 
  男子和女子一英里赛的最佳拟合直线是否相交？
  \begin{answer}
    男子和女子一英里赛直线方程为
    （女子一英里赛方程的纵截距项已根据上面的答案调整，使其在年份 $0$ 处为零，因为男子一英里赛方程采用了同样处理。）
    \begin{align*}
        y &=994.8276974-0.3871993827x  \\
        y &=3125.6426-1.445753225x
    \end{align*}
    显然两直线相交。
    使用计算机程序最容易完成
    运算：MuPAD 给出 $x=2012.949004$ 和 $y=215.4150856$
    ($215$~seconds is $3$~minutes and $35$~seconds).
    \textit{Remark.}
    当然，这种投影存在很大疑问 \Dash  例如，
    女子方程受早期较慢纪录的影响
    \Dash  但过程仍然很有趣。
     \begin{center}  \small
       \includegraphics{map/mp/ch3.89}
     \end{center}
  \end{answer}
%  \item 
%     Look up the records for the men's and women's marathon.
%     求出 line of best fit for each.
%     Do the lines cross?
%  \item 
%    Look up the records for men's long jump.
%    How far out-of-trend was Bob Beamon's record?
  \item
    \textit{（这说明某些数据集并不适合线性模型，
    此时最佳拟合直线
    不具有预测价值。）}
    在公路餐馆，一名卡车司机告诉我，他的老板经常让他绕路， 
    虽然耗油更多
    但桥梁通行费更低。
    他说纽约州会校准
    哈德逊河每座桥的通行费，
    权衡从纽约市前往该处的额外燃油费
    与较低的过桥费，
    以鼓励人们前往州北部。
    下表引自 \cite{CostOfTolls} 和 \cite{GoogleMaps}，
    列出哈德逊河各收费 crossing 的
    从时代广场驾车的英里数及美元费用
    （乘用车；若单向收费则列出其一半）。
    (if a crossings has a one-way toll then it shows half that number). 
    \begin{center}
      \small
      % \begin{tabular}{lrr}
      %    \multicolumn{1}{c}{\textit{Crossing}}
      %      &\multicolumn{1}{c}{\textit{Distance}}
      %      &\textit{Toll}  \\ \hline
      %    Lincoln Tunnel &$2$  &$6.00$ \\
      %    Holland Tunnel &$7$  &$6.00$  \\
      %    George Washington Bridge &$8$  &$6.00$  \\
      %    Verrazano-Narrows Bridge &$16$  &$6.50$  \\
      %    Tappan Zee Bridge &$27$  &$2.50$  \\
      %    Bear Mountain Bridge &$47$  &$1.00$  \\
      %    Newburgh-Beacon Bridge &$67$  &$1.00$  \\
      %    Mid-Hudson Bridge &$82$  &$1.00$  \\
      %    Kingston-Rhinecliff Bridge &$102$  &$1.00$  \\
      %    Rip Van Winkle Bridge &$120$  &$1.00$  
      % \end{tabular}
      \begin{tabular}{ccc}
         \textit{Crossing}
           &\textit{Distance}
           &\textit{Toll}  \\ \hline
         \begin{tabular}{@{}l@{}}
           Lincoln Tunnel \\
           Holland Tunnel  \\
           George Washington Bridge  \\
           Verrazano-Narrows Bridge   \\
           Tappan Zee Bridge  \\
           Bear Mountain Bridge  \\
           Newburgh-Beacon Bridge  \\
           Mid-Hudson Bridge   \\
           Kingston-Rhinecliff Bridge  \\
           Rip Van Winkle Bridge   
         \end{tabular}
         &\begin{tabular}{@{}r@{}}
           $2$ \\
           $7$  \\
           $8$  \\
           $16$  \\
           $27$   \\
           $47$   \\
           $67$   \\
           $82$  \\
           $102$  \\
           $120$ 
          \end{tabular}  
         &\begin{tabular}{@{}r@{}}
           $6.00$ \\
           $6.00$  \\
           $6.00$  \\
           $6.50$  \\
           $2.50$  \\
           $1.00$  \\
           $1.00$  \\
           $1.00$  \\
           $1.00$  \\
           $1.00$ 
         \end{tabular} 
      \end{tabular}
    \end{center}
    求出最佳拟合直线并绘图，以
    说明司机是在利用我的轻信。
    \begin{answer}
      \textit{Sage} gives the line of best fit as
      $\text{toll}=-0.05\cdot\text{dist}+5.63$.
\begin{lstlisting}
sage: dist = [2, 7, 8, 16, 27, 47, 67, 82, 102, 120]
sage: toll = [6, 6, 6, 6.5, 2.5, 1, 1, 1, 1, 1]
sage: var('a,b,t')
(a, b, t)
sage: model(t) = a*t+b
sage: data = zip(dist,toll)
sage: fit = find_fit(data, model, solution_dict=True)
sage: model.subs(fit)
t |--> -0.0508568169130319*t + 5.630955848442933
sage: p = plot(model.subs(fit), (t,0,120))+points(data,size=25,color='red')
sage: p.save('bridges.pdf')        
\end{lstlisting}
      但图表显示该方程几乎没有预测价值。
      \begin{center}
        \includegraphics[width=.5\textwidth]{map/pix/bridges.pdf}
      \end{center}
      显然更好的模型是：
      （只有一个中间例外）市区的 crossing
      费用大致相同，
      州北部 crossing 的费用也彼此相同。
    \end{answer}
  \item 
    1986 年挑战者号航天飞机爆炸时，对 NASA 发射决定的批评之一是其
    对
    O 形密封圈失效次数与温度关系的分析方式
    （O 形密封圈失效导致爆炸）。
    四次 O 形密封圈失效将造成灾难性后果。
    NASA 掌握此前 $24$ 次飞行的数据。
    \begin{center}
      \begin{tabular}{r|rrrrrrrrrrr}
         \textit{temp\ ${}^\circ$F} 
            &53 &75 &57 &58 &63 &70 &70 &66 &67 &67 &67 \\
         \hline
         \textit{failures}  
            &3  &2  &1  &1  &1  &1  &1  &0  &0  &0  &0
      \end{tabular}\hspace*{3em}                                           \\
      \hspace*{3em}\begin{tabular}{|rrrrrrrrrrrrr}
                68 &69 &70 &70 &72 &73 &75 &76 &76 &78 &79 &80 &81\\
         \hline
                0  &0  &0  &0  &0  &0  &0  &0  &0  &0  &0  &0  &0
      \end{tabular}
    \end{center}
    当天预报温度为 \(31^\circ\text{F}\)。
    \begin{exparts}
      \partsitem NASA based the decision to launch partially on a chart showing
        仅包含至少发生一次 O 形密封圈失效的飞行。
        求出最适合这七次飞行的直线。
        根据这些数据，
        预测温度为 $31$ 时 O 形密封圈失效次数，并预测
        失效次数超过四次时的温度。
      \partsitem 求出最适合全部 24 次飞行的直线。
        根据这些额外数据，
        预测温度为 $31$ 时 O 形密封圈失效次数，并预测
        失效次数超过四次时的温度。
    \end{exparts}
    你认为哪种预测方法更准确？
    （优秀的讨论见 \cite{Stats}。）
    \begin{answer}
      \begin{exparts}
        \partsitem MAPLE 或 MuPAD 等计算机代数系统给出
          截距为 $b=4259/1398\approx 3.239628$
          斜率为 $m=-71/2796\approx -0.025393419$     
          将 $x=31$ 代入方程，预测
          O 形密封圈失效次数为 $y=2.45$（四舍五入到两位）。
          代入 $y=4$ 并求解，得到温度
          $x=-29.94^\circ$F.
        \partsitem On the basis of this information
          \begin{equation*}
            A =  
            \begin{mat}
               1 & 53 \\
               1 & 75 \\
%              1 & 57 \\
%              1 & 58 \\
%              1 & 63 \\
%              1 & 70 \\
%              1 & 70 \\
%              1 & 66 \\
%              1 & 67 \\
%              1 & 67 \\
%              1 & 67 \\
%              1 & 68 \\
%              1 & 69 \\
%              1 & 70 \\
%              1 & 70 \\
%              1 & 72 \\
%              1 & 73 \\
%              1 & 75 \\
%              1 & 76 \\
%              1 & 76 \\
%              1 & 78 \\
%              1 & 79 \\
               \vdots \\
               1 & 80 \\
               1 & 81
            \end{mat}
            \qquad
            b = 
            \colvec{ 3 \\
                     2 \\
%                    1 \\
%                    1 \\
%                    1 \\
%                    1 \\
%                    1 \\
%                    0 \\ 
%                    0 \\
%                    0 \\
%                    0 \\
%                    0 \\
%                    0 \\
%                    0 \\
%                    0 \\
%                    0 \\
%                    0 \\
%                    0 \\
%                    0 \\
%                    0 \\
%                    0 \\
%                    0 \\
                     \vdots \\
                     0 \\
                     0}
          \end{equation*}
          MAPLE 给出截距 $b=187/40=4.675$，
          slope $m=-73/1200\approx -0.060833$.
          此时将 $x=31$ 代入方程预测
          $y=2.79$ O-ring failures (rounded to two places).
          代入 $y=4$ 次失效得到温度
          $x=11^\circ$F。
     \begin{center}  \small
       \includegraphics{map/mp/ch3.90}
      \end{center}
      \end{exparts}  
    \end{answer}
  \item 
     下表列出太阳到前七颗行星的平均距离，
     以地球平均距离为单位。
     \begin{center}
       \begin{tabular}{ccccccc}
         Mercury &Venus   &Earth   &Mars    &Jupiter &Saturn  &Uranus  \\ 
         \hline     
         $0.39$  &$0.72$  &$1.00$  &$1.52$  &$5.20$  &$9.54$  &$19.2$
       \end{tabular}
     \end{center}
    \begin{exparts}
      \partsitem 绘制行星编号
        （水星为 \(1\)，依此类推）与距离的关系图。
        注意图形不像直线，因此求
        最佳拟合直线没有意义。
      \partsitem 但它看起来像指数曲线。因此
        绘制行星编号与
        距离对数的关系图。
        这看起来像直线吗？
      \partsitem 火星与木星之间的小行星带是
        
        一颗解体行星的残余。
        重新编号，使木星为 $6$、土星为 $7$、天王星为
        $8$，再次绘制与对数的关系图。
        这次看起来更好吗？
      \partsitem 对这些数据使用最小二乘法预测
       海王星的位置。
      \partsitem 重复上述过程预测冥王星的位置。
      \partsitem 该公式对海王星和冥王星准确吗？ 
    \end{exparts}
    这种方法曾帮助发现海王星（尽管第二小题关于历史的说法
    并不准确；实际上，海王星在
    位于 $9$ 号位置
    $9$ 号位置的发现促使人们寻找位于 $5$ 号位置的“缺失行星”。）
    See \cite{Gardner70}
      \begin{answer}
        \begin{exparts}
          \partsitem The plot is nonlinear.
            \begin{center}  \small
              \includegraphics{map/mp/ch3.91}
            \end{center}
          \partsitem Here is the plot.
           \begin{center}  \small
             \includegraphics{map/mp/ch3.92}
           \end{center}
           行星 $4$ 与 $5$ 之间可能有向上的突变。
          \partsitem This plot seems even more linear.
           \begin{center}  \small
             \includegraphics{map/mp/ch3.93}
           \end{center}
          \partsitem
            使用以下输入
            \begin{equation*}
              A =  
              \begin{mat}[r]
                1 & 1 \\
                1 & 2 \\
                1 & 3 \\
                1 & 4 \\
                1 & 6 \\
                1 & 7 \\
                1 & 8
              \end{mat}
              \qquad
              b = \colvec[r]{-0.40893539 \\
                          -0.1426675 \\
                           0 \\
                           0.18184359 \\
                           0.71600334 \\
                           0.97954837 \\
                           1.2833012}
            \end{equation*}
            MuPAD 给出截距 $b= -0.6780677466$，
            斜率为 $m=0.2372763818$。
          \partsitem Plugging $x=9$ into the equation
            $y= -0.6780677466+0.2372763818x$ from the prior item gives
            距离对数为 $1.4574197$，因此预测
            距离为 $28.669472$。
            实际距离约为 $30.003$。
          \partsitem Plugging $x=10$ into the same equation
            给出距离对数为 $1.6946961$，因此预测
            距离为 $49.510362$。
            实际距离约为 $39.503$。
        \end{exparts}
      \end{answer}
  % \item 
  %   William Bennett has proposed an Index of Leading Cultural Indicators
  %   for the US (\cite{Bennett}, in 1993).
  %   Among the statistics cited are the average daily hours spent watching TV, 
  %   and the average combined SAT scores.
  %   \begin{center}
  %     \begin{tabular}{r|cccccccc}
  %        &1960   &1965   &1970   &1975   &1980   &1985   &1990   &1992  
  %           \\ \hline
  %        \textit{TV}
  %        &5:06   &5:29   &5:56   &6:07   &6:36   &7:07   &6:55   &7:04 \\
  %        \textit{SAT}
  %        &$975$  &$969$  &$948$  &$910$  &$890$  &$906$  &$900$  &$899$ 
  %     \end{tabular} 
  %   \end{center}
  %   设 a cause and effect relationship is proposed between
  %   the time spent watching TV and the decline in SAT scores (in this article,
  %   Mr.~Bennett does not argue that there is a direct connection).
  %   \begin{exparts}
  %    \partsitem 求出 line of best fit relating the independent variable
  %      of average daily TV hours to the dependent variable of SAT scores.
  %    \partsitem 求出 most recent estimate of the average daily TV hours  
  %      (Bennett's cites Neilsen Media Research as the source of these 
  %      estimates).
  %      Estimate the associated SAT score.
  %      How close is your estimate to the actual average?
  %      (Warning: a change has been made recently in the SAT,
  %      so you should investigate whether some adjustment needs to be made to
  %      the reported average to make a valid comparison.)
  %   \end{exparts}
  %   \begin{answer}
  %     \begin{exparts}
  %       \partsitem With this input
  %         \begin{equation*}
  %           A = 
  %           \begin{mat}[r]
  %             1 & 306 \\
  %             1 & 329 \\
  %             1 & 356 \\
  %             1 & 367 \\
  %             1 & 396 \\
  %             1 & 427 \\
  %             1 & 415 \\
  %             1 & 424
  %           \end{mat}
  %           \qquad
  %           b = 
  %           \colvec[r]{975 \\
  %                   969 \\
  %                   948 \\
  %                   910 \\
  %                   890 \\
  %                   906 \\
  %                   900 \\
  %                   899}
  %         \end{equation*}
  %         MAPLE gives the intercept 
  %         $b=34009779/28796\approx 1181.0591$ and the slope 
  %         $m=-19561/28796\approx -0.6793$.
  %         \begin{center}  \small
  %           \includegraphics{map/mp/ch3.94}
  %      \end{center}
  %     \end{exparts}
  %   \end{answer}
%  \item
%    Derive the 
%    \definend{normal equations}\index{least squares!normal equations}%
%    \index{normal equations!for least squares} for least squares line fitting:
%    for data points \( (x_1,y_1) \), \ldots, \( (x_n,y_n) \)
%    the line of best fit \( y=mx+b \) has this slope
%    \begin{equation*}
%      m=\frac{\frac{\displaystyle x_1y_1+\dots+x_ny_n}{\displaystyle n}
%              -(\frac{\displaystyle x_1+\dots+x_n}{\displaystyle n})
%               (\frac{\displaystyle y_1+\dots+y_n}{\displaystyle n})}{
%              \frac{\displaystyle {x_1}^2+\cdots+{x_n}^2}{\displaystyle n}
%              -(\frac{\displaystyle x_1+\dots+x_n}{\displaystyle n})^2}
%    \end{equation*}
%    and this vertical intercept.
%    \begin{equation*}
%      b=\frac{y_1+\dots+y_n}{n}-m(\frac{x_1+\dots+x_n}{n})
%    \end{equation*}
%    (\textit{Remark}.
%     A advantage of these equations is that a
%     person doing by-hand calculations can make up \Dash  and check \Dash  a 
%     table with a
%     column of $x_i$'s, one of $y_i$'s, one of $x_i^2$'s,
%     one of $y_i^2$'s , and finally one of $x_i\cdot y_i$'s.
%     Then the sums of these columns fit into the formulas.)
%     \begin{answer}
%       The $\nbyn{2}$ shows how the larger cases hold.
%     \end{answer}
\item 设 $W$ 是某个 $\Re^n$ 的子空间，并且
  $\vec{v}$ is not an element of~$W$.
  设 $\vec{v}$ 到 $W$ 的正交投影为向量
  $\proj{\vec{v}}{W}=\vec{p}$.
  证明 $\vec{p}$ 是 $W$ 中距离 $\vec{v}$ 最近的元素。
  \begin{answer}
    For any $\vec{w}\in W$, the vectors $\vec{v}-\vec{p}$ and
    $\vec{p}-\vec{w}$ are orthogonal.
    因此三角不等式可应用于以这些向量为边的三角形，
    两边，且以 $\vec{v}-\vec{w}$ 为斜边。
    Therefore $\vec{v}-\vec{w}$ is at least as long as $\vec{v}-\vec{p}$.
  \end{answer}

   





\end{exercises}
%\par\noindent\announcecomputercode %
%\begin{computercode}
%    #!/usr/bin/python
%    # least_squares.py   calculate the line of best fit for a data set
%    # data file format: each line is two numbers, x and y
%    import math, string
%
%    n = 0
%    sum__x = 0
%    sum_y = 0
%    sum_x_squared = 0
%    sum_xy = 0
%
%    fn = raw_input("Name of the data file? ")
%    datafile = open(fn,"r")
%    while 1:
%      ln = datafile.readline()
%      if (ln):
%            data = string.split(ln)
%            x = string.atof(data[0])
%            y = string.atof(data[1])
%            n = n+1
%            sum_x = sum_x+x
%            sum_y = sum_y+y
%            sum_x_squared = sum_x_squared+(x*x)
%            sum_xy = sum_xy +(x*y)
%      else:
%            break
%    datafile.close()
%
%    slope=(sum_of_xy-pow(sum_x,2)/n)/(sum_x_squared-pow(sum_x,2)/n)
%    intercept=(sum_y-(slope*sum_x))/n
%    print "line of best fit: slope=",slope," intercept=",intercept 
%\end{computercode}


\index{line of best fit|)}
\index{least squares|)}





